% 2026-11-11 / v36
\section{線束生成部分の飽和と全coherent残余圏}
\label{crs:sec:saturation}

本節は、一般点復元で用いた線束生成部分と全coherent導来圏の差を解消する。
鍵は、一般点で既に得た線束生成残余圏のsmoothnessを、
全coherent対象に対する有限性と組み合わせることである。
特殊点の生成性を一般点へ持ち上げる議論は使わない。

既知の入力は、Burke--Stevensonの導来section、Gorenstein環の
$\chi$条件と有限cohomological dimension、およびOrlov型分解
\cite[Proposition~4.18, Lemmas~6.1--6.2, Theorem~6.4]{CRSBS}、
smooth dg圏上のpseudo-perfect加群のperfect性
\cite[Lemma~2.8(2)]{CRSTV}、admissibleな半直交分解とsmoothness
\cite[Lemma~4.15(3)]{CRSPerry}である。
これらの一般原理自体の新規性は主張しない。
本研究での寄与は、その仮定を一般点の中心三次商について確認し、
線束生成残余pairの復元を全coherent残余pairへ拡張する組合せにある。
この組合せの包括的新規性調査・外部査読は未実施である。

\subsection{抽象的な飽和判定}

$k$を体、$\mathcal D$を小さな冪等完備pretriangulated $k$-線形dg圏とし、
$G\in\mathcal D$、$\mathcal C=\operatorname{thick}(G)$、
$\Lambda=\operatorname{RHom}_{\mathcal D}(G,G)$と置く。
以下の直交は全cohomological shiftに関するものである。

\begin{theorem}[有限性とsmoothnessによる全体への飽和]
\label{crs:thm:abstract}
次の三条件を仮定する。
\begin{enumerate}
\item $\Lambda$は$k$上smooth、すなわち
$\Delta_\Lambda\in\operatorname{Perf}(\Lambda^{\mathrm{op}}\otimes_k\Lambda)$である。
\item 任意の$M\in\mathcal D$について
$\operatorname{RHom}(G,M)\in\operatorname{Perf}(k)$である。
\item $\mathcal C^\perp=0$である。
\end{enumerate}
このとき$\mathcal C=\mathcal D$であり、
$\operatorname{Perf}(\Lambda)\to\mathcal D$はdg Morita同値である。
\end{theorem}
\begin{proof}
右$\Lambda$-加群$N_M=\operatorname{RHom}(G,M)$を考える。
smoothな対角の有限cell/retract表示へ
$N_M\otimes_\Lambda^{\mathbf L}(-)$を適用する。
自由な包絡加群の像は$N_M\otimes_k\Lambda$であり、
$N_M\in\operatorname{Perf}(k)$だから右$\Lambda$-perfectである。
従って$N_M\in\operatorname{Perf}(\Lambda)$となる。
これは\cite[Lemma~2.8(2)]{CRSTV}の一対象版の直接証明である。

$G$を用いるMorita関手
$F:\operatorname{Perf}(\Lambda)\xrightarrow{\sim}\mathcal C$
で$N_M$を移し、評価射$F(N_M)\to M$を取る。
ここで$F(N_M)$は有限coneと直和因子だけで作られ、
$\mathcal D$に無限直和があることを要求しない。
$F$と$\operatorname{RHom}(G,-)$の単位は$\Lambda$で同型なので、
有限操作による帰納法で全perfect加群に対して同型である。
従って評価射に$\operatorname{RHom}(G,-)$を施すと同型となる。
評価射のconeは$G$、従って$\mathcal C$と右直交する。
第三条件によりこのconeは零である。
全$M$が$\mathcal C$に属し、enhancedな完全忠実性は
$\operatorname{thick}(G)$におけるMorita表示から従う。
\end{proof}

\begin{corollary}[対角の有限表示からの一様な生成長]
\label{crs:cor:level}
上の対角が自由包絡加群から$q$段の拡大、有限直和、shift、直和因子で作れるなら、
任意の$M\in\mathcal D$も$G$から高々$q$段の同じ操作で作れる。
この系は$q$の存在から生成長を評価するもので、具体的な値は与えない。
\end{corollary}
\begin{proof}
前証明の各自由包絡加群の像は$N_M\otimes_k\Lambda$である。
体上では$N_M$はその有限次元有界cohomologyの直和に準同型だから、
この像は有限個の$\Lambda$のshiftの直和である。
テンソル関手と$F$は拡大段数を増やさない。
\end{proof}

\begin{example}[smoothnessを省けない]
$\Lambda=k[\epsilon]/(\epsilon^2)$、
$\mathcal D=D^b(\operatorname{mod}\Lambda)$、$G=\Lambda$とする。
$\operatorname{RHom}(G,M)$は全$M$について$k$-perfectで、$G^\perp=0$である。
しかし$k$の射影分解は$\epsilon$による無限周期分解であり、
$\operatorname{Tor}^{\Lambda}_i(k,k)=k$が全$i\ge0$に残る。
従って$\operatorname{thick}(G)=\operatorname{Perf}(\Lambda)$は$\mathcal D$より真に小さい。
単なるright-spanning性を有限生成性と同一視してはいけない。
\end{example}

\subsection{中心三次商のcoherent対象に対する有限性と検出}

$A$を連結局所有限な両側noetherian Koszul AS-regular $k$-代数とし、
global dimensionとGorenstein parameterをともに$N\ge5$とする。
$w\in Z(A)_3$を正則元、$B=A/(w)$、$a=N-3$と置く。
$\operatorname{qgr}B$は有限生成次数付き右加群を有限次元加群で割った圏、
$\mathcal O(j)=\pi B(j)$とする。標準のenhancementを用いる。

\begin{lemma}[必要なGorenstein・section条件]
\label{crs:lem:sections}
$B$は両側noetherian AS-Gorensteinで、parameterは$a$である。
$B$は$\chi$条件と有限cohomological dimensionを満たす。
任意の$M\in D^b(\operatorname{qgr}B)$と任意の$j\in\mathbb Z$について
\[
 \operatorname{RHom}_{\operatorname{qgr}B}(\mathcal O(j),M)
 \in\operatorname{Perf}(k).
 \tag{CS1}\label{crs:eq:finite-Hom}
\]
また、全$j,p$で$\operatorname{Hom}(\mathcal O(j),M[p])=0$なら$M=0$である。
\end{lemma}
\begin{proof}
$0\to A(-3)\xrightarrow{w}A\to B\to0$と導来change-of-ringsにより、
$B$-加群$M$について
\[
 R\underline{\operatorname{Hom}}_B(M,B)
 \simeq R\underline{\operatorname{Hom}}_A(M,A)(-3)[1]
\]
を得る。ここでは$w$の中央性を使い、normal元のtwistを省略したのではない。
左右でinjective dimensionは有限であり、$M=k$を代入すると
$R\underline{\operatorname{Hom}}_B(k,B)\simeq k(N-3)[-(N-1)]$となる。
$B$のnoetherian性は商であることから従う。
従って\cite[Lemmas~6.1--6.2]{CRSBS}を適用できる。

\cite[Proposition~4.18とその直前の導来随伴]{CRSBS}により、任意の整数$b$について
\[
 R\Gamma_{\ge b}:D^b(\operatorname{qgr}B)
 \longrightarrow D^b(\operatorname{gr}_{\ge b}B)
 \tag{CS2}\label{crs:eq:sections}
\]
は完全忠実で、各cohomologyは有限生成次数付き$B$-加群である。
固定内部次数の成分は有限次元であり、cohomologyは有界だから、
$\bigl(R\Gamma_*M\bigr)_{-j}=R\operatorname{Hom}(\mathcal O(j),M)$は
$k$-perfectとなる。任意の$j$に対し$b\le-j$を選べばよい。

全ての線束とのHomが零なら$R\Gamma_{\ge0}M$の各次数のcohomologyが零である。
この複体自体が零となり、完全忠実性から$M=0$を得る。
この検出は高次Extも含み、射の次数零だけを調べているのではない。
\end{proof}

\begin{lemma}[有限線束生成対象]
\label{crs:lem:finite-lines}
\[
 \mathcal C_B=\operatorname{thick}\{\mathcal O(j):j\in\mathbb Z\},
 \qquad G_B=\bigoplus_{j=0}^{N-1}\mathcal O(j)
\]
とすると$\mathcal C_B=\operatorname{thick}(G_B)$であり、$\mathcal C_B$はproperである。
\end{lemma}
\begin{proof}
$k$の$A$上の有限Koszul分解を$B$へ導来テンソルする。
二項の$A$-分解で計算すると$\operatorname{Tor}^A_i(k,B)$は
$i=0,1$にのみあり、それぞれ$k,k(-3)$である。
従ってこのKoszul複体を$qgr$へ移すと完全になる。
両端は一つの線束で、中間の項のshiftは$1,\ldots,N-1$にある。
任意の連続$N$個の線束から、一つ手前と一つ後をそれぞれ生成できる。
前後への帰納により全てのtwistを得る。
proper性は補題\ref{crs:lem:sections}を有限個の$\mathcal O(i),\mathcal O(j)$に適用し、
有限coneと直和因子へ延長して従う。
\end{proof}

\begin{theorem}[線束部分のsmoothnessで全coherent圏を得る]
\label{crs:thm:coherent}
上の中心三次商について、$\mathcal C_B$がsmoothなら
\[
 \boxed{\quad
 \mathcal C_B=D^b(\operatorname{qgr}B).
 \quad}\tag{CS3}\label{crs:eq:full}
\]
従って右辺はsmooth proper dg圏で、$G_B$が古典的生成対象になる。
\end{theorem}
\begin{proof}
任意のcoherent有界対象に対する$\operatorname{RHom}(G_B,-)$の有限性は
補題\ref{crs:lem:sections}による。線束による検出と
補題\ref{crs:lem:finite-lines}から$\mathcal C_B^\perp=0$である。
定理\ref{crs:thm:abstract}を適用する。
ここで全$D^b(\operatorname{qgr}B)$のsmoothnessを仮定していない。
\end{proof}

\subsection{残余部分のsmoothnessと有限拡大からの降下}

\begin{lemma}[有限体拡張からのsmoothness降下]
\label{crs:lem:descent}
有限体拡大$k'/k$と$k$-dg代数$\Gamma$について、
$\Gamma\otimes_k k'$が$k'$上smoothなら$\Gamma$は$k$上smoothである。
\end{lemma}
\begin{proof}
$E=\Gamma^{\mathrm{op}}\otimes_k\Gamma$と置く。
$E\otimes_k k'$は$E$上有限自由なので、そのperfect加群のrestrictionは$E$-perfectである。
従ってsmoothな対角$\Gamma\otimes_k k'$をrestrictionしたものは$E$-perfect。
$\lambda:k'\to k$を$\lambda(1)=1$となる$k$-線形写像とすれば、
$\Gamma\to\Gamma\otimes_k k'\xrightarrow{1\otimes\lambda}\Gamma$は
$E$-線形なretractである。従って$\Gamma$は$E$-perfectである。
\end{proof}

\begin{proposition}[残余のsmoothnessから全線束部分へ]
\label{crs:prop:residual-smooth}
$\mathcal K_B^{\rm lin}=\mathcal C_B\cap
\langle\mathcal O,\ldots,\mathcal O(a-1)\rangle^\perp$とする。
$\mathcal K_B^{\rm lin}$がsmoothなら$\mathcal C_B$もsmoothである。
従ってこの場合
\[
 \mathcal K_B^{\rm lin}=\mathcal K_B^{\rm coh}:=
 D^b(\operatorname{qgr}B)\cap
 \langle\mathcal O,\ldots,\mathcal O(a-1)\rangle^\perp.
 \tag{CS4}\label{crs:eq:residual-full}
\]
\end{proposition}
\begin{proof}
$\mathcal O,\ldots,\mathcal O(a-1)$は例外列である。
これはGorensteinの線束cohomology公式、または
\cite[Theorem~6.4]{CRSBS}のOrlov分解から従う。
$\mathcal C_B$はproperなので、例外列への評価・余評価は有限であり、
その生成部分はadmissibleである。残余部分もadmissibleとなり
\[
 \mathcal C_B=\langle\mathcal K_B^{\rm lin},
                 \mathcal O,\ldots,\mathcal O(a-1)\rangle
\]
を得る。smoothnessの貼合せ\cite[Lemma~4.15(3)]{CRSPerry}を適用する。
具体的には各例外成分は$\operatorname{Perf}(k)$でsmooth、
相互Homはproperであり、smoothな成分の間の貼合せ双加群は
\cite[Lemma~2.8(2)]{CRSTV}によりperfectである。
上三角dg圏の対角の有限三角を順に貼り合わせれば全体の対角もperfectとなる。
定理\ref{crs:thm:coherent}から\eqref{crs:eq:residual-full}を得る。
\end{proof}

\subsection{一般点復元を全coherent残余pairへ拡張する}

\begin{theorem}[全coherent残余pairの幾何学的一般点復元]
\label{crs:thm:generic}
$K=\mathbb C(h)$とし、既存の一般点復元定理の入力\textup{(L),(CF)}を満たす
一変数量子化$B_K=A_K/(\Omega_K)$を考える。
$A_K$は上のnoetherian Koszul AS-regular仮定を満たすとする。
この仮定は本研究の局所量子化入力の一般ファイバーで成立する。
既存定理が与える有限拡大$K'/K$と滑らかな三次式$W'/K'$について
\[
 \boxed{\quad
 (\mathcal K^{\rm coh}_{B_{K'}},\rho_B)
 \simeq(\operatorname{Ku}(V(W')),\rho_{W'}).
 \quad}\tag{CS5}\label{crs:eq:generic}
\]
さらに$D^b(\operatorname{qgr}B_K)=\mathcal C_{B_K}$であり、
このdg圏と$\mathcal K^{\rm coh}_{B_K}$は$K$上smooth properである。
混合$\lambda=1$で既存入力\textup{(D)}も用いる場合、$W'$は既に計算した
多項式$G_h+v^3$に取れる。
\end{theorem}
\begin{proof}
既存の次数シフト軌道による一般点定理を用いると、
$\mathcal K^{\rm lin}_{B_{K'}}$は滑らかな可換cubicのKuと同値でありsmooth properである。
線束生成子のHomの平坦底変更も入力(L)に含まれているため、
残余生成代数について$\Gamma_{K'}=\Gamma_K\otimes_KK'$と同定できる。
補題\ref{crs:lem:descent}で$\Gamma_K$のsmoothnessが従う。
命題\ref{crs:prop:residual-smooth}と定理\ref{crs:thm:coherent}を$K$上で適用する。

同じ議論は$K'$上でも成立する。有限拡大後のnoetherian性は
$A_K\otimes_KK'$が$A_K$上有限加群であることから従う。
標数零なので$K'/K$は分離的であり、任意の加群は
restriction後のscalar extensionの直和因子になる。
この分離idempotentと$A_K$上の有限射影分解によりglobal dimensionの上界を保つ。
Koszul分解とGorensteinの等式は平坦底変更で保たれる。
従って\eqref{crs:eq:residual-full}を$K'$にも適用でき、
既存の線束部分の同値がそのまま全coherent残余pairの同値となる。
rotationは同じ$L_{\mathcal O}\circ(1)$であり、新たな関手への変更はない。
(D)による既存の明示式同定を合成すると最後の結論を得る。
\end{proof}

\begin{remark}[主張の境界]
本定理は局所window比較(L)と完備MF比較(CF)を用いる。
これらは定理\ref{tw:thm:window}および定理\ref{cft:thm:complete}で構成されている。
特定の可換方程式を指定する場合には、さらに比較(D)を用いる。
全$D^b(\operatorname{qgr}B_{K'})$と$D^b(\operatorname{coh}V(W'))$の同値は主張しない。
残余成分が同値でも例外成分との貼合せは異なり得る。
またアフィン次数付き環$B_K$のglobal dimensionが有限とは言っていない。
特殊点を含む同値のspread、同値核の行列・有限拡大次数、
複数の基底チャートにわたる相対同値の接着は別の問題である。
\end{remark}

\subsection{既知の生成原理と適用条件の確認}

StevensonのTheorem~6.13\cite{CRSStevenson}は、affineな局所超曲面環について
\emph{全ての}thick部分圏を分類する。tensor idealに限定した定理ではない。
本研究で使う$K[x]/(W)$は原点のみが特異なので、適用範囲は合っている。
次数の再付与では有限対象から無限直和へ移るが、有限rank MFからの射は
有限個の直和成分を通るため、最後のcompact局所化によって有限thick構成へ戻る。
従って、次数の再付与に現れる無限直和は最終的な生成性の結論ではなく、
有限thick閉包に戻るための中間的な構成である。

完備MF同値の局所化では、先に固定有限対象へ制限してから$h$を反転する。
各固定cohomological次数は有限自由な斉次成分からなり、
平坦局所化はそのHomとcohomologyに可換する。
任意のLaurent係数対象へ新しく無限Taylor和を評価することは使わない。
曲率付き関手と局所windowの構成は、それぞれ
第\ref{cft:sec:transport}節と第\ref{tw:sec:window}節に与える。

全coherent残余圏への拡張には、上で示したsection有限性、
smooth対角、評価射および直交成分の検出を用いる。
有限係数のEuler計算は非零性を調べるものであって、
これらの圏論的議論を置き換えるものではない。
